\section{The \texorpdfstring{$RO(C_{p^n})$}{RO(Cpn)}-graded Tambara functor}
\label{sec:uniform}
Let $R$ be a nonzero $p$-torsionfree perfectoid ring. We compute the
$RO(C_{p^n})$-graded Tambara functor of $X_R$.
The homotopy groups of the local terms are computed mainly
by combining earlier results. We give these computations in
Appendix~\ref{app:spectral-foundations}.
We first compute norms in natural degrees.
We then compute norms on the homotopy groups of $X_R^{t_i}$.
Finally, we use the boundary maps of \eqref{eq:tate-compression}
to compute norms in shifted degrees.

\subsection{The Breuil--Kisin twist}\label{subsec:period-line}
We need the following
Breuil--Kisin twist to describe the homotopy groups.

\begin{definition}\label{def:period-line}
Let $R$ be a nonzero $p$-torsionfree perfectoid ring with perfect prism
$(A,I_0,\varphi)$. Thus $A=W(R^\flat)$, where
$R^\flat=\varprojlim_{x\mapsto x^p}R/pR$, $I_0=\ker\theta$ for the
Fontaine map $\theta:A\to R$, and $\varphi$ is Witt Frobenius.
Put $I_j=\varphi^j(I_0)$. For $0\le i\le m$, write
$J_i^m=\prod_{i<j\le m}I_j$, with $J_i^i=A$. Define
\begin{equation*}
B=\operatorname{BK}(A,I_1)
   =\varprojlim_{r\ge1}\left(J_0^r/(J_0^r)^2,\,\tau_r\right).
\end{equation*}
The transition map $\tau_r:J_0^{r+1}/(J_0^{r+1})^2\to J_0^r/(J_0^r)^2$
is uniquely determined by
\[
p\tau_r([x])=[x]\qquad(x\in J_0^{r+1}).
\]
Existence and uniqueness follow from
\cite[Corollary~2.2.10]{BhattLurie}.
\end{definition}

By \cite[Constructions 2.2.11, 2.2.14 and
Proposition 2.2.12]{BhattLurie}, $B$ is an invertible $A$-module and
\begin{equation}\label{eq:algebraic-period-phi}
 \Phi_B:\varphi^*B\xrightarrow{\sim}I_1^{-1}B,
 \qquad \varphi(J_0^r)=I_1^{-1}J_0^{r+1}.
\end{equation}
Write
\[
\operatorname{can}:
\operatorname{TC}^-(R;\Z_p)=X_R^{h\mathbb T}
\longrightarrow X_R^{t\mathbb T}=\operatorname{TP}(R;\Z_p).
\]
Use the complex orientation under
$\mathbb T/C_{p^r}\xrightarrow{\sim}\mathbb T$,
$[z]\mapsto z^{p^r}$.
By \cite[Proposition~6.5 and Remark~6.6]{BMS},
\[
\begin{tikzcd}[column sep=large,row sep=small]
\pi_2\operatorname{TC}^-(R;\Z_p)
  \arrow[r,"\pi_2(\operatorname{can})"] \arrow[d,"\cong"']
& \pi_2\operatorname{TP}(R;\Z_p) \arrow[d,"\cong"] \\
I_0B \arrow[r,hook] & B.
\end{tikzcd}
\]
For $r\ge1$, the cyclotomic map
$X_R^{hC_{p^{r-1}}}\to X_R^{tC_{p^r}}$
is $\mathbb T$-equivariant.
In the homotopy fixed-point spectral sequences for the residual circle,
use its complex Euler class to identify $\pi_2$ with the first
quotient. Naturality \cite[Remark~6.6]{BMS} gives
\[
\begin{tikzcd}[column sep=large,row sep=large]
\pi_2 X_R^{hC_{p^{r-1}}}
  \arrow[r] \arrow[d,"\cong"']
& \pi_2 X_R^{tC_{p^r}} \arrow[d,"\cong"] \\
\dfrac{I_0J_0^{r-1}}{(I_0J_0^{r-1})^2}
  \arrow[r,"\varphi"']
& \dfrac{J_0^r}{(J_0^r)^2}.
\end{tikzcd}
\]
The inverse limit and \cite[Construction~2.2.14]{BhattLurie}
for $(A,I_1)$ give the degree-two linearization of cyclotomic Frobenius:
\[
\varphi^*(I_0B)=I_1\varphi^*B
 \xrightarrow{\ I_1\Phi_B\ }B.
\]
We use the canonical isomorphism
\begin{equation}\label{eq:period-cotangent-reduction}
B/I_0B\xrightarrow{\sim}R,\qquad
I_0B/I_0^2B\xrightarrow{\sim}I_0/I_0^2
\end{equation}
of \cite[Proposition~6.5 and the preceding discussion]{BMS}. By \cite[Theorem~6.1]{BMS},
$\pi_{2a}X_R\cong(I_0/I_0^2)^a$ for $a\ge0$. Projection to $r=1$ and iteration of \eqref{eq:algebraic-period-phi} give
\begin{equation}\label{eq:period-repair-data}
B/I_1B\simeq I_1/I_1^2,\qquad
 \Sigma_r:\varphi^{r*}B\xrightarrow{\sim}(J_0^r)^{-1}B\quad(r\ge0).
\end{equation}
When $I_0$ is principal, we can use the first isomorphism
in \eqref{eq:period-repair-data} to choose a generator of $B$.

\begin{proposition}\label{lem:period-repair-frame}
If $I_0=(\xi)$, every lift $b\in B$ of
$\varphi(\xi)\bmod I_1^2$ generates $B$ freely over $A$.
\end{proposition}
\begin{proof}
The ring $A$ is $(p,I_1)$-complete, so $I_1\subseteq\operatorname{Jac}(A)$.
The map $A\to B$, $1\mapsto b$, is an isomorphism modulo $I_1$.
By Nakayama's lemma it is surjective. Since $B$ is invertible, it is
an isomorphism.
\end{proof}

\subsection{Homotopy groups in representation degrees}
We first reduce to $RO_{\mathrm{sel}}(C_{p^n})$
using maps of representation spheres.

\begin{definition}\label{def:character-translation}
For $0\le j<n$, $r>0$, and $(r,p)=1$, the map
\[
f_{r,j}:S^{\lambda_{p^n}(p^j)}
\longrightarrow S^{\lambda_{p^n}(rp^j)},
\qquad f_{r,j}(z)=z^r,
\]
is a $p$-local equivalence. Its map on $C_{p^l}$-fixed spheres has
degree $r$ for $l\le j$ and degree $1$ for $l>j$. Define the unit
\[
\widetilde v_{r,j}=r f_{r,j}^{-1}
\in
\pi_{\lambda_{p^n}(rp^j)-\lambda_{p^n}(p^j)}
^{C_{p^n}}\mathbb S_{(p)},
\]
with fixed-point degrees
\[
\deg(\widetilde v_{r,j}^{C_{p^l}})
=
\begin{cases}
1,&l\le j,\\
r,&l>j.
\end{cases}
\]
For
\[
D=\sum_{r,j}m_{r,j}
\bigl(\lambda_{p^n}(rp^j)-\lambda_{p^n}(p^j)\bigr),
\qquad m_{r,j}\in\mathbb Z,
\]
define
\[
\widetilde v_D
=\prod_{r,j}\widetilde v_{r,j}^{m_{r,j}}.
\]
The unit $\widetilde v_D$ depends on the weights $r$. Under restriction we use the same maps $f_{r,j}$.
Let $v_{r,j},v_D$ be the images of $\widetilde v_{r,j},\widetilde v_D$ in $X_R$.
For $V=V_0+D$ with $V_0$ selected, multiplication
by $v_D$ induces an isomorphism
\[
\pi_{V_0}^{C_{p^n}}X_R
\xrightarrow{\ \cong\ }
\pi_V^{C_{p^n}}X_R.
\]
\end{definition}

For $V\in RO(C_{p^n})$, obtain $V_0$ by replacing
$\lambda_{p^n}(rp^j)$ by $\lambda_j$ in its irreducible
expansion, where $0<rp^j<p^n/2$ and $(r,p)=1$,
and leaving the one-dimensional summands unchanged.
Use the expansion of $D=V-V_0$ in
Definition~\ref{def:character-translation}.
The local homotopy groups in degrees $V_0$ and $V$ are
identified by multiplication by $v_D$,
which commutes with $\ell_i$ and $\rho_i$.

We have the following definition.

\begin{definition}\label{def:natural-degree}\label{def:operation-complexes}
Let $R$ be a nonzero $p$-torsionfree perfectoid ring.
Fix $n\ge0$ and $V\in\RO(C_{p^n})$, and put
\[
U=\begin{cases}V,&\dim V^{C_{p^n}}\equiv0\pmod2,\\
V+1,&\dim V^{C_{p^n}}\equiv1\pmod2,
   \end{cases}
\qquad a_i=\frac12\dim_{\mathbb R}U^{C_{p^i}}\quad(0\le i\le n).
\]
For $p$-torsionfree $R$, define $A_i(U)$ and $T_i(U)$ for $0\le i<n$
as follows. Products of invertible ideals and their
negative powers are taken in the localization of $A$ obtained by
inverting their generators.

If $\epsilon(\mathbf a)=0$, set
\begin{align}
E_i(U)&=(\varphi^{n*}B)^{a_n}\prod_{i<j\le n}I_j^{a_j}
\quad(0\le i\le n),\notag\\
A_i(U)&=\frac{I_i^{a_i^+}E_i(U)}
                {J_i^n I_i^{(a_i+1)^+}E_i(U)},
 &T_i(U)&=E_i(U)/J_i^n E_i(U).\label{eq:uniform-complex}
\end{align}

If $\epsilon(\mathbf a)=1$, set
\begin{align}
\mathcal S_i(U)&=(\varphi^{n*}B)^{a_n}I_n^{a_n}\prod_{i<j<n}I_j^{a_j+1/2},\notag\\
A_i(U)&=I_i^{(a_i+1/2)^+}\mathcal S_i(U)/I_nI_i^{(a_i+1/2)^+}\mathcal S_i(U),
 &T_i(U)&=\mathcal S_i(U)/I_n\mathcal S_i(U).\label{eq:uniform-sign}
\end{align}
For $p$-torsionfree $R$, set
\[
A_n(U)=\begin{cases}
\displaystyle
\frac{I_n^{a_n}(\varphi^{n*}B)^{a_n}}
     {I_n^{a_n+1}(\varphi^{n*}B)^{a_n}},&a_n\ge0,\\[6pt]
0,&a_n<0.
\end{cases}
\]
For $0\le i<n$, define $\ell_i:A_i(U)\to T_i(U)$ and
$\rho_i:A_{i+1}(U)\to T_i(U)$ as follows.
For $p$-torsionfree $R$, both maps send $[x]$ to $[x]$, with
representatives in $E_i(U)$ if $\epsilon(\mathbf a)=0$ and in
$\mathcal S_i(U)$ if $\epsilon(\mathbf a)=1$.
For $0\le i<n-1$, use the identifications
\[
E_i(U)=I_{i+1}^{a_{i+1}}E_{i+1}(U),
\qquad
\mathcal S_i(U)=I_{i+1}^{a_{i+1}+1/2}\mathcal S_{i+1}(U).
\]
For $n\ge1$ and $a_n\ge0$, we have $A_n(U)=T_{n-1}(U)$ and
$\rho_{n-1}=\mathrm{id}$; for $a_n<0$, the source is zero.

For $G=C_{p^n}$, we use the incidence map to define the following complex.
\[
\begin{gathered}
C_G^\bullet(U)=
\left[\bigoplus_{i=0}^n A_i(U)\xrightarrow{\Delta_U}
\bigoplus_{i=0}^{n-1}T_i(U)\right],\\
(\Delta_Uq)_i=\ell_i(q_i)-\rho_i(q_{i+1}),
\end{gathered}
\]
with cochain degrees $0,1$.
\end{definition}

Sulyma computes finite-level $\TR$ groups in a
contiguous range of representation degrees~\cite[Theorem~4.12]{SulymaNew}.
For every natural degree $U$, the kernel and cokernel of $\Delta_U$ give
the homotopy groups in degrees $U$ and $U-1$, respectively.

\begin{proposition}\label{thm:uniform}
For $R$, $n$ and $U$ as in Definition~\ref{def:natural-degree},
\begin{equation}\label{eq:uniform-groups}
 \pi_U^{C_{p^n}}X_R=\ker\Delta_U,
 \qquad \pi_{U-1}^{C_{p^n}}X_R=\coker\Delta_U.
\end{equation}
\end{proposition}

\begin{proof}
For $U\in RO_{\mathrm{sel}}(C_{p^n})$, the identifications
\[
A_i(U)=\pi_U^{C_{p^n}}X_R^{h_i},\qquad
T_i(U)=\pi_U^{C_{p^n}}X_R^{t_i},\qquad
A_n(U)=\pi_U^{C_{p^n}}X_R^{\Phi_n}
\]
with the first two taken for $0\le i<n$, follow from
Proposition~\ref{prop:app-finite-comparison},
\eqref{eq:app-frob-numerator}--\eqref{eq:app-frob-denominator}, and
\eqref{eq:app-sign-local-kernels}.
Multiplication by $v_D$ gives the same identifications
for all natural degrees. The same computations give
\[
\pi_{U\pm1}^{C_{p^n}}X_R^{h_i}=0,\qquad
\pi_{U\pm1}^{C_{p^n}}X_R^{t_i}=0,\qquad
\pi_{U\pm1}^{C_{p^n}}X_R^{\Phi_n}=0.
\]
The long exact sequence of the fiber of the incidence map
in Section~\ref{sec:compression} gives \eqref{eq:uniform-groups}.
\end{proof}

\subsection{Products, restrictions, transfers, and Weyl actions}
\label{subsec:products-maps}

We now give the formulas for products, restrictions, transfers, and Weyl actions.

When comparing $G=C_{p^n}$ with $H=C_{p^{n-1}}$, insert a zero
summand at $A_n$ and $T_{n-1}$ in $C_H^\bullet$.

For natural degrees $U,W$ with degree coordinates
$\mathbf a,\mathbf b$ satisfying $\epsilon(\mathbf a)=\epsilon(\mathbf b)=0$,
the products at $A_i$ and $T_i$ for $i<n$ are induced by
$E_i(U)\otimes_A E_i(W)\simeq E_i(U+W)$.

At $p=2$, when $\epsilon(\mathbf a)=1$, we use the sign Euler class to
express restriction and transfer as connecting homomorphisms.

\begin{proposition}\label{prop:sign-exact-complexes}
Let $G=C_{2^n}$, $H=C_{2^{n-1}}$, and $n\ge1$.
If $U$ is natural with $\epsilon(\mathbf a)=1$,
there are degreewise exact sequences
\begin{equation}\label{eq:sign-exact-complexes}
\begin{aligned}
0\longrightarrow C_H^\bullet(\Res_H^G(U+\sigma_{2^n}))
&\xrightarrow{\tr}C_G^\bullet(U+\sigma_{2^n})
\xrightarrow{a_{\sigma_{2^n}}}C_G^\bullet(U)\longrightarrow0,\\
0\longrightarrow C_G^\bullet(U)
&\xrightarrow{a_{\sigma_{2^n}}}C_G^\bullet(U-\sigma_{2^n})
\xrightarrow{\res}C_H^\bullet(\Res_H^G(U-1))\longrightarrow0.
\end{aligned}
\end{equation}
Both maps $a_{\sigma_{2^n}}$ send $[q]$ to $[q]$.
The connecting homomorphisms on cohomology are
\[
\begin{aligned}
\res_H^G&:H^0(C_G^\bullet(U))
\longrightarrow H^1(C_H^\bullet(\Res_H^G(U+\sigma_{2^n}))),\\
\tr_H^G&:H^0(C_H^\bullet(\Res_H^G(U-1)))
\longrightarrow H^1(C_G^\bullet(U)).
\end{aligned}
\]
\end{proposition}

\begin{proof}
Firstly we apply $(G/H)_+\to S^0\to S^{\sigma_{2^n}}$ and its dual to
\eqref{eq:tate-compression}. By Shapiro's equivalence and the vanishing
in shifted degrees, we have  \eqref{eq:sign-exact-complexes}.

Write $\sigma=\sigma_{2^n}$. For $i<n$, the coefficient maps are
\[
A_i(U+\sigma_{2^n})\twoheadrightarrow A_i(U)
\hookrightarrow A_i(U-\sigma_{2^n}),
\]
and likewise for $T_i$.
Indeed, $E_i(U+\sigma)=\mathcal S_i(U)=J_i^{n-1}E_i(U-\sigma)$.
The second map is the kernel inclusion in \eqref{eq:app-sign-local-kernels}.
Their composite is $a_{2\sigma}=a_{\lambda_{n-1}}$, which acts by inclusion
on the circle representatives \cite[Theorem~1.9 and Example~3.2]{SulymaNew}.
Hence the first map sends $[q]$ to $[q]$ as well. At $A_n$, $\sigma^G=0$.

By naturality,  the connecting homomorphisms are $\res_H^G$
and $\tr_H^G$.
\end{proof}

For $U=\rho_p(\mathbf a)$ with $\mathbf a\in D_n^{(0)}$, we use the same virtual complex
$\mathbb T$-representation lifting $U$ to compute the terms at
$G$ and at $H$.
For $V=V_0+D$, we can transport the sphere maps by
\begin{equation}\label{eq:character-res-tr}
\begin{aligned}
\res_H^G(v_Dx_0)
 &=\res_H^G(v_D)\res_H^G(x_0),\\
\tr_H^G(y)
 &=v_D\,\tr_H^G\!\left(\res_H^G(v_D)^{-1}y\right).
\end{aligned}
\end{equation}
Here $x_0\in\pi_{V_0}^G X_R$ and
$y\in\pi_{\Res_H^G V}^H X_R$.
The above identities follow from naturality and the projection formula.
For $p$-torsionfree $R$, \cite[Proposition~2.2.12]{BhattLurie}
for $(A,I_1)$ and \eqref{eq:algebraic-period-phi} give
\[
 B/(I_0J_0^r)B\simeq J_0^r/(I_0J_0^r)J_0^r,\qquad
 B/J_0^rB\simeq J_0^r/(J_0^r)^2\quad(r\ge0).
\]
The first follows from $B/J_0^{r+1}B\simeq J_0^{r+1}/(J_0^{r+1})^2$
by pulling back along $\varphi^{-1}$ and tensoring with $I_0^{-1}$.
For $U=\rho_p(\mathbf a)$ with $\mathbf a\in D_n^{(0)}$, the Frobenius isomorphism gives
\[
E_i(U)\cong
(\varphi^{(n-1)*}B)^{a_n-a_{n-1}}\otimes_A E_i(\Res_H^G U),
\qquad 0\le i<n,
\]
where the two lines are formed at $G$ and $H$, respectively.
The quotient isomorphisms and $\varphi^{i*}\Sigma_{n-1-i}$ identify
$\varphi^{(n-1)*}B$ with $A$ modulo $I_iJ_i^{n-1}$.
Thus, using the same circle lift, the terms at $H$ are
\begin{equation}\label{eq:subgroup-coefficients}
\begin{aligned}
A_i(\Res_H^G U)&\cong
\frac{I_i^{a_i^+}E_i(U)}
     {J_i^{n-1}I_i^{(a_i+1)^+}E_i(U)}
&& (0\le i<n),\\
T_i(\Res_H^G U)&\cong E_i(U)/J_i^{n-1}E_i(U)
&& (0\le i<n-1).
\end{aligned}
\end{equation}

We now define the transfer coefficient.

\begin{definition}\label{def:transfer-coefficient}
For $m\ge0$, let
$V:W_{m+1}(R)\to W_{m+2}(R)$
be the Witt vector Verschiebung.
For $p$-torsionfree $R$, the maps
$\widetilde\theta_r:A\to W_r(R)$
of \cite[Lemma~3.12]{BMS18} give
\[
\widetilde\theta_{m+1}\varphi:
A/(I_0J_0^m)\xrightarrow{\sim}W_{m+1}(R),
\qquad
v_m=(\widetilde\theta_{m+2}\varphi)^{-1}(V(1))
\in A/(I_0J_0^{m+1}).
\]
\end{definition}

We can now give the formulas.

\begin{proposition}\label{prop:products-maps}\label{lem:path-products}
Let $R$ be as in Definition~\ref{def:natural-degree} and
let $G=C_{p^n}$.

\begin{enumerate}[label=\textup{(\arabic*)},leftmargin=*]
\item For natural degrees $U,W$ with degree coordinates
$\mathbf a,\mathbf b$, respectively, let
\[
x\in\ker\Delta_U,\quad y\in\ker\Delta_W,\quad
[z]\in\coker\Delta_U,\quad [w]\in\coker\Delta_W.
\]
The unit in degree zero is $(1,\ldots,1)$.
Under \eqref{eq:uniform-groups}, the products are
\begin{equation}\label{eq:path-cup}
\begin{aligned}
(xy)_j&=x_jy_j &&(0\le j\le n),\\
x[w]&=[(\ell_i(x_i)w_i)_{i<n}],\\
[z][w]&=0\in\pi_{U+W-2}^G X_R.
\end{aligned}
\end{equation}

\item The restriction and transfer maps on kernels and cokernels are as follows.
Let $n\ge1$ and $H=C_{p^{n-1}}$. For a natural degree $U$
with $\epsilon(\mathbf a)=0$, the maps
\[
\begin{aligned}
\res_H^G&:C_G^\bullet(U)\longrightarrow C_H^\bullet(\Res_H^G U),\\
\tr_H^G&:C_H^\bullet(\Res_H^G U)\longrightarrow C_G^\bullet(U)
\end{aligned}
\]
are given, for $p$-torsionfree $R$ and selected $U$, in the
coordinates of \eqref{eq:subgroup-coefficients}, by
\begin{equation}\label{eq:adjacent-res-tr}
\begin{aligned}
\res_H^G(x_0,\ldots,x_n)&=( x_0,\ldots, x_{n-1}),\\
\tr_H^G(y_0,\ldots,y_{n-1})&=(v_{n-1}y_0,\ldots,v_{n-1}y_{n-1},0),\\
\res_H^G(z_0,\ldots,z_{n-1})&=( z_0,\ldots, z_{n-2}),\\
\tr_H^G(w_0,\ldots,w_{n-2})&=(v_{n-1}w_0,\ldots,v_{n-1}w_{n-2},0).
\end{aligned}
\end{equation}
For arbitrary $U$, use \eqref{eq:character-res-tr}.

\item Let $p=2$, $n\ge1$, and $H=C_{2^{n-1}}$.
Suppose $U$ is natural with
$\epsilon(\mathbf a)=1$. For
\[
x\in\ker\Delta_U,\qquad
y\in\ker\Delta_{\Res_H^G(U-1)},
\]
choose lifts
\[
b\in C_G^0(U+\sigma_{2^n}),\quad a_{\sigma_{2^n}} b=x,
\qquad
c\in C_G^0(U-\sigma_{2^n}),\quad\res_H^G(c)=y.
\]
Then
\begin{equation}\label{eq:sign-res-tr}
\begin{aligned}
\res_H^G(x)&=[\tr^{-1}\Delta_{U+\sigma_{2^n}}b]
\in\coker\Delta_{\Res_H^G(U+\sigma_{2^n})},\\
\tr_H^G(y)&=[a_{\sigma_{2^n}}^{-1}\Delta_{U-\sigma_{2^n}}c]
\in\coker\Delta_U.
\end{aligned}
\end{equation}
Moreover,
\[
\tr_H^G\big|_{\pi_{\Res_H^G U}^H X_R}=0,
\qquad
\res_H^G\big|_{\pi_{U-1}^G X_R}=0.
\]
\end{enumerate}
\end{proposition}

\begin{proof}
\textup{(1)} The maps in diagram of \eqref{eq:tate-compression}
are maps of commutative $X_R$-algebras in $\Sp^G$, so the
associated fiber sequence is a sequence of $X_R$-modules.
By the proof of \cite[Theorem~4.30, equations~(4.31)--(4.32)]{ZhangHZ} and \eqref{eq:uniform-groups}, we have the three product
formulas in \eqref{eq:path-cup}.

For $p$-torsionfree $R$, tensoring the ideals in
\eqref{eq:uniform-complex} and \eqref{eq:uniform-sign}
gives the products formula.
For $\epsilon(\mathbf a)=\epsilon(\mathbf b)=1$ and $0\le i<n$,
multiplication is induced by the inclusion
\[
\mathcal S_i(U)\mathcal S_i(W)
=J_i^{n-1}E_i(U+W)\subseteq E_i(U+W).
\]
The unit has coordinate $1$ at every vertex.

\textup{(2)} Restriction and transfer commute with the incidence map
and its connecting homomorphism. Thus they induce the maps on
kernels and cokernels, as in the proof of
\cite[Theorem~6.13]{ZhangHZ}.
For $p$-torsionfree $R$,
\cite[Proposition~6.2 and Remark~6.6]{BMS} and
\cite[Theorem~3.3]{HM97} give
\[
\tr_H^G(1)
=(\widetilde\theta_{n+1}\varphi)^{-1}(V(1))
=v_{n-1}.
\]
Restriction is the quotient in \eqref{eq:subgroup-coefficients} at each vertex and edge.
Transfer follows from the projection formula, proving \eqref{eq:adjacent-res-tr}.

\textup{(3)} The proofs of
\cite[Propositions~6.17--6.18]{ZhangHZ}, applied to
\eqref{eq:sign-exact-complexes}, give
\eqref{eq:sign-res-tr} up to sign.
Since $X_R$ has a circle action and $\epsilon(\mathbf a)=1$,
\[
\begin{aligned}
\res_H^G(x)
 &=c_{\gamma_{2^n}}\res_H^G(x)=-\res_H^G(x),\\
\tr_H^G(y)
 &=\tr_H^G(c_{\gamma_{2^n}}y)=-\tr_H^G(y).
\end{aligned}
\]
The proof of \cite[Proposition~6.20]{ZhangHZ} gives
$\tr_H^G|_{\pi_{\Res_H^G U}^H X_R}=0$ and
$\res_H^G|_{\pi_{U-1}^G X_R}=0$:
the map
\[
X_R\longrightarrow
\left(\prod_{i<n}X_R^{h_i}\right)\times X_R^{\Phi_n}
\]
induces zero on the source groups and an injection on the target groups
by Proposition~\ref{thm:uniform}.
\end{proof}

It remains to determine how a generator of $G$ acts on
$\pi_{\Res_L^G V}^L X_R$ for $L<G$.

\begin{proposition}
\label{prop:weyl-action}
Let $G=C_{p^n}$, $V\in RO(G)$, and let $L<G$ be any proper
subgroup. Let $\gamma$ be the generator of $G$.
For odd $p$, the Weyl action on
$\pi_{\Res_L^G V}^L X_R$ is trivial.
For $p=2$, write
\[
V=m+s\sigma_{2^n}+\sum_\lambda s_\lambda\lambda,
\]
where $\lambda$ ranges over the irreducible real rotation
representations. Then
\[
c_{\gamma^a}\big|_{\pi_{\Res_L^G V}^L X_R}
=(-1)^{as}\operatorname{id}
\qquad(a\in\mathbb Z).
\]
\end{proposition}

\begin{proof}
For $f\in[S^{\Res_L^G V},X_R]^L$,
\[
c_\gamma(f)
=\gamma_{X_R}\circ f\circ\gamma_{S^V}^{-1}.
\]
A path in $\mathbb T$ from $1$ to $\gamma$ gives
$\gamma_{X_R}\simeq_L\operatorname{id}$.
For $p=2$,
\[
L\leq\ker\sigma_{2^n},\qquad
\gamma_{S^{\sigma_{2^n}}}\simeq_L-\operatorname{id},
\qquad
\gamma_{S^\lambda}\simeq_L\operatorname{id},
\]
where the last homotopy is through rotations.
The sphere identities extend to virtual summands, so
\[
\gamma_{S^V}\simeq_L(-1)^s\operatorname{id},
\qquad
c_{\gamma^a}(f)=(-1)^{as}f.
\]
For odd $p$, every nontrivial irreducible real representation
is a rotation, so we have $c_\gamma=\operatorname{id}$.
\end{proof}

\subsection{Norms}\label{subsec:graded-norms}
We now compute the norms.

Let $H=C_{p^h}<G=C_{p^n}$. We first compute
\[
N_H^G:\pi_V^H X_R\longrightarrow\pi_{\Ind_H^G V}^G X_R,
\qquad V\in RO_{\mathrm{sel}}(H).
\]

\subsubsection{Fixed-point coordinates of induced representations}\label{subsubsec:selected-degree-norms}
For $V\in RO(H)$, write
\[
\begin{aligned}
2a_j&=\dim V^{C_{p^j}} &&(0\le j\le h),\\
2A_j&=\dim(\Ind_H^G V)^{C_{p^j}} &&(0\le j\le n).
\end{aligned}
\]

To determine the target degree of $N_H^G$, we compute the $A_j$
from the $a_j$.

\begin{proposition}\label{prop:induction-of-degrees}
The degree coordinates of $\Ind_H^G V$ satisfy
\begin{equation}\label{eq:period-repair-induced-degree}
A_j=\begin{cases}
p^{n-h}a_j,&0\le j\le h,\\
p^{n-j}a_h,&h<j\le n.
\end{cases}
\end{equation}
In particular, for $n=h+1$,
\begin{equation*}
I(\mathbf a)=(pa_0,\ldots,pa_h,a_h).
\end{equation*}
If $V$ is natural, then $\Ind_H^G V$ is natural and
$\epsilon(\mathbf A)=0$.
\end{proposition}

\begin{proof}
For $K=C_{p^j}$, the fixed subspace of the induced representation is
\[
(\Ind_H^G V)^K\cong
\bigoplus_{K\backslash G/H}V^{K\cap H}.
\]
Here $K\cap H=C_{p^{\min(j,h)}}$ and
$|K\backslash G/H|=p^{n-\max(j,h)}$, which gives
\eqref{eq:period-repair-induced-degree}. The same formula holds
for virtual representations by additivity.
\end{proof}

\subsubsection{Norms in natural degrees}
Let $U\in RO_{\mathrm{sel}}(H)$ be natural.
We use geometric fixed points to determine $N_H^G$ on $\pi_U^H X_R$.

We first give the following definition.

\begin{definition}\label{def:geometric-fixed-point-maps}
For $U\in RO(C_{p^n})$, define
\[
\begin{aligned}
\gamma_j\colon
\pi_U^{C_{p^n}}X_R
&\xrightarrow{\Phi^{C_{p^j}}}
\pi_{U^{C_{p^j}}}^{C_{p^{n-j}}}
\Phi^{C_{p^j}}X_R\\
&\xrightarrow[\cong]{\mathrm{cyclotomic}}
\pi_{U^{C_{p^j}}}^{C_{p^{n-j}}}X_R
\xrightarrow{\operatorname{res}_e^{C_{p^{n-j}}}}
\pi_{U^{C_{p^j}}}^{e}X_R,
\qquad 0\le j\le n.
\end{aligned}
\]
For \(m\ge 1\), let
\[
\begin{gathered}
[m]\colon S^1\longrightarrow S^1,
\qquad z\longmapsto z^m,\\
\psi^m\colon X_R\longrightarrow [m]^*X_R,
\end{gathered}
\]
in $\Sp^{B\mathbb T}$, where $[m]^*$ pulls back the circle action along $[m]$ and
$\psi^m$ is the Adams operation induced by $[m]$
\cite[Definition~5.1 and Remark~5.2]{SulymaNew}.
\end{definition}

For $p$-torsionfree $R$, the Adams operations satisfy the following.
On $\pi_*X_R$, we have $\psi^{p^r}=(\psi^p)^r$.
By \cite[Remark~5.2 and Lemma~5.3]{SulymaNew},
\begin{equation}\label{eq:rcyc-crossing}
R_{\mathrm{cyc}}N_e^{C_p}=\psi^p,
\qquad
\left.\psi^{p^r}\right|_{\pi_{2d}X_R}
=p^{rd}\operatorname{id}
\quad(r,d\ge0).
\end{equation}
The first identity uses pullback along $\varrho_p$ and,
for a finite-dimensional real vector space $W$, the identification
\[
W\xrightarrow{\sim}(\Ind_e^{C_p}W)^{C_p},
\qquad w\longmapsto(w,\ldots,w).
\]
For negative summands, use the duals of the inverse sphere maps.

To recover a norm from its images under the maps $\gamma_j$, we need
to know when these maps detect a class.

\begin{lemma}\label{lem:geometric-ghost-injectivity}
Let $R$ be $p$-torsionfree and let $U\in RO(H)$ be natural, with
degree coordinates $\mathbf a=(a_0,\ldots,a_h)$. The maps
\[
\begin{aligned}
(\gamma_0,\ldots,\gamma_h):\pi_U^H X_R
&\longrightarrow\prod_{j=0}^h\pi_{U^{C_{p^j}}}^eX_R,
&&\epsilon(\mathbf a)=0,\\
\gamma_h:\pi_U^H X_R&\longrightarrow\pi_{U^H}^eX_R,
&&\epsilon(\mathbf a)=1,
\end{aligned}
\]
are injective. 
\end{lemma}
\begin{proof}
Multiplication by $v_D$ reduces the assertion to
$U\in RO_{\mathrm{sel}}(H)$.
By \cite[Lemma~2.2.8]{BhattLurie},
$\theta\varphi^{-i}(I_j)=pR$ for $j>i$.
For an invertible $A$-module $E$, a local generator of $I_j$
therefore acts injectively on $E/I_iE$, since $R$ is $p$-torsionfree.
The $I_i$-adic filtration gives injectivity on $E/I_i^bE$ for every
$b\ge0$, also for a local generator of $J_i^h$. Hence
\begin{equation}\label{eq:ghost-intersection}
\begin{aligned}
I_i^bE\cap I_jE&=I_i^bI_jE,\qquad 0\le i<j\le h,\\
I_i^bE\cap J_i^hE&=I_i^bJ_i^hE,\qquad 0\le i<h.
\end{aligned}
\end{equation}

Suppose $\epsilon(\mathbf a)=0$ and
$x=(x_0,\ldots,x_h)\in\ker\Delta_U$ satisfies $\gamma_jx=0$ for
$0\le j\le h$. Cyclotomic equivalence gives $x_h=0$.
If $x_{i+1}=0$, then
\[
\ell_i(x_i)=\rho_i(x_{i+1})=0.
\]
For $a_i\ge0$, \eqref{eq:ghost-line-conversion} and $\gamma_ix=0$
place a representative of $x_i$ in $I_i^{a_i+1}E_i(U)$.
Since $\ell_i(x_i)=0$, this representative also lies in $J_i^hE_i(U)$.
Equation~\eqref{eq:ghost-intersection} gives $x_i=0$.
For $a_i<0$, $\ell_i$ is an isomorphism.
Induction gives $x=0$.

Suppose $\epsilon(\mathbf a)=1$.
Equation~\eqref{eq:ghost-intersection} with $j=h$ implies that
the maps $\ell_i$ of \eqref{eq:uniform-sign} are injective for $0\le i<h$.
If $\gamma_hx=0$, then $x_h=0$ and
$\ell_i(x_i)=\rho_i(x_{i+1})$ gives $x_i=0$ for $i=h-1,\ldots,0$.

\end{proof}

Let $M$ be the ordered model of $\Ind_H^G U$ and
$\eta:S^M\xrightarrow{\sim}S^{\Ind_H^G U}$ its grading map
\cite[Definition~5.1 and Remark~5.2]{GradedTambara}.
For each $\chi:G\to\{\pm1\}$, by $\sqrt{[G:H]}$ times evaluation at $H$
we identify the $\chi$-isotypic subspace of $\Ind_H^G U$ with the ordered
$\chi|_H$-isotypic subspace of $U$. For $K=C_{p^j}$ let
\[
d_j:\bigoplus_{O\in K\backslash G/H} U^{K\cap H}
  \xrightarrow{\sim}(\Ind_H^G U)^K
\]
be the orbit diagonal. Order orbits by their least cosets
in the order $H,\gamma_{p^n}H,\ldots$, and set
$d_j(v)=|O|^{-1/2}v$ at the least coset of $O$.

\begin{proposition}\label{prop:vertex-norms}
Let $R$ be $p$-torsionfree and let $U\in RO_{\mathrm{sel}}(H)$ be natural.
Write $M=M_0+D$ as in Definition~\ref{def:character-translation}.
For $x\in\pi_U^H X_R$, the element
$y=v_D^{-1}\eta^*N_H^Gx\in\ker\Delta_{M_0}$ is given by
\begin{equation*}
\gamma_j(y)=\frac{\deg(d_j^{-1}\Phi^{C_{p^j}}\eta)}{\gamma_j(v_D)}
\begin{cases}
(\gamma_jx)^{p^{n-h}},&0\le j\le h,\\
\bigl(\psi^{p^{j-h}}\gamma_hx\bigr)^{p^{n-j}},&h<j\le n.
\end{cases}
\end{equation*}
\end{proposition}

\begin{proof}
In the orbit basis $d_j$, $\Phi^K$ gives $[G:HK]$ factors
\cite[Appendix~A]{HHR}, each equal to $\gamma_jx$ for $j\le h$
and $\psi^{p^{j-h}}\gamma_hx$ for $j>h$
by \eqref{eq:rcyc-crossing}. Pullback along $\Phi^K\eta$ contributes
$\deg(d_j^{-1}\Phi^K\eta)$. Uniqueness follows from
Proposition~\ref{prop:induction-of-degrees} and
Lemma~\ref{lem:geometric-ghost-injectivity}.
\end{proof}

In degree zero this recovers the Witt norm
\cite[Theorems~1.4 and~4.1]{WittNorm} and its prismatic form
\cite[Section~3.1]{SulymaPrisms}.

For products and composites, go from the induced grading map
$\eta'$ to $\eta$ acts by the sphere unit $(\eta')^{-1}\eta$.

\subsubsection{Norms on $X_R^{t_i}$}
For $0\le i<h<n$, Proposition~\ref{prop:boundary-indexed-multiplications} constructs
compatible $G/H$-indexed multiplications on the spectra in the
Tate compression diagram truncated at $X_R^{\Phi_{i+1}}$.
For each spectrum $Y$ in this diagram, write
\[
\mu_Y:N_H^G\res_H^G Y\longrightarrow Y.
\]
For each arrow $f:Y\to Z$, we have
\[
f\mu_Y=\mu_ZN_H^G(\res_H^Gf).
\]
For $U\in RO(H)$, write
\[
N_H^G:\pi_U^H X_R^{t_i}
\longrightarrow\pi_{\Ind_H^G U}^G X_R^{t_i}
\]
for the norm induced by $\mu_{X_R^{t_i}}$.
Write $\partial_i:\Omega X_R^{t_i}\to X_R$ and
$\partial_{i*}:\pi_U^K X_R^{t_i}\to\pi_{U-1}^K X_R$ for the $i$th boundary
of \eqref{eq:tate-compression} and its induced map, where
$K\in\{H,G\}$ and $U\in RO(K)$.

Let $R$ be $p$-torsionfree and $0\le i<h$.
Put $\lambda_j=\lambda_{p^h}(p^j)$ for $0\le j<h$.

We have the following definition.

\begin{definition}\label{def:intrinsic-orientations}
For $0\le j<h$, define
\[
\begin{aligned}
\mathcal O_{j,h}:(J_0^j)^{-1}B&\longrightarrow
\pi_{2-\lambda_j}^H X_R=\ker\Delta_{2-\lambda_j},\\
\ell&\longmapsto([\ell]_r)_{0\le r\le h}.
\end{aligned}
\]
Here $[\ell]_r$ is the class of $\ell$ in $A_r(2-\lambda_j)$.
Using $\Sigma_h$, the modules of representatives are
$(J_0^j)^{-1}B$ for $r\le j$ and $(J_0^{r-1})^{-1}B$ for $r>j$.
The inclusions $(J_0^j)^{-1}B\subseteq (J_0^{r-1})^{-1}B$ give
$\ell_r([\ell]_r)=\rho_r([\ell]_{r+1})$.
Under $\pi_\star^H X_R\to\pi_\star^H X_R^{t_i}$, a class
$\mathcal O_{j,h}(b')$ with $b'$ an $A$-basis element of $(J_0^j)^{-1}B$
maps to a unit for $i\le j<h$. The Euler class $a_{\lambda_j}$ maps to
a unit for $0\le j\le i$.

Choose $I_0=(\xi)$ and $b$ as in
Proposition~\ref{lem:period-repair-frame}. Put
\[
u_{\lambda_j}=\mathcal O_{j,h}\left(
\frac{b}{\prod_{l=1}^j\varphi^l(\xi)}\right).
\]
For $U=\rho_H(\mathbf a)$, $\mathbf a\in D_h^{(0)}$, the $A/J_i^h$-module
$\pi_U^H X_R^{t_i}=T_i(U)$ is free of rank one with basis element
\begin{equation}\label{eq:tate-unit-generator}
g_{T_i}^H(\mathbf a)=
\left(\frac{u_{\lambda_i}}{a_{\lambda_i}}\right)^{a_{i+1}}
\prod_{i<j<h}u_{\lambda_j}^{a_{j+1}-a_j}\prod_{0\le j\le i}a_{\lambda_j}^{a_{j+1}-a_j}.
\end{equation}
\end{definition}

Under $\Sigma_h^{a_h}$, this basis is represented by
\[
b^{a_h}\prod_{1\le r\le i}\varphi^r(\xi)^{-a_h}
\prod_{i<r\le h}\varphi^r(\xi)^{a_r-a_h}.
\]

For $s\in A/J_i^h$, choose a lift
$\widetilde s\in\pi_0^H X_R=A/(I_0J_0^h)$. Naturality gives
\[
N_H^G(s)=[N_H^G(\widetilde s)]\in A/J_i^n,
\]
where Proposition~\ref{prop:vertex-norms} with $U=0$ computes the
right-hand side. The result is independent of the lift.

It remains to compute the norm of the basis element $g_{T_i}^H(\mathbf a)$.

\begin{proposition}\label{lem:period-repair-norm}
Let $R$ be $p$-torsionfree, $0\le i<h$, and
$U=\rho_H(\mathbf a)$, $\mathbf a\in D_h^{(0)}$.
For $s\in A/J_i^h$, with $g_{T_i}^H(\mathbf a)$ as in
Definition~\ref{def:intrinsic-orientations}, we have
\[
\begin{aligned}
N_H^G(sg_{T_i}^H(\mathbf a))
={}&N_H^G(s)
\left(\frac{N_H^G(u_{\lambda_i})}
{a_{\Ind_H^G\lambda_i}}\right)^{a_{i+1}}\\
&\cdot\prod_{i<j<h}
N_H^G(u_{\lambda_j})^{a_{j+1}-a_j}
\prod_{0\le j\le i}
a_{\Ind_H^G\lambda_j}^{a_{j+1}-a_j}
\in T_i(\Ind_H^G U).
\end{aligned}
\]
\end{proposition}

\begin{proof}
Apply multiplicativity and
$N_H^G(a_{\lambda_j})=a_{\Ind_H^G\lambda_j}$
to \eqref{eq:tate-unit-generator}.
Use Proposition~\ref{prop:vertex-norms} for $N_H^G(u_{\lambda_j})$. Thus the proof follows.
\end{proof}

\subsubsection{Norms in shifted degrees}
Let $R$ be $p$-torsionfree and $U=\rho_H(\mathbf a)$, $\mathbf a\in D_h$.
Put $V=U-1$ and $T=\Ind_H^G V+1$.
By Proposition~\ref{prop:induction-of-degrees}, the degree coordinates of $T$ are
\begin{equation*}
\tfrac12\dim T^{C_{p^j}}=
\begin{cases}
p^{n-h}a_j-(p^{n-h}-1)/2,&0\le j\le h,\\
p^{n-j}a_h-(p^{n-j}-1)/2,&h<j\le n.
\end{cases}
\end{equation*}
For $n=h+1$, put
\begin{equation*}
J(\mathbf a)=(pa_0-(p-1)/2,\ldots,pa_h-(p-1)/2,a_h).
\end{equation*}

In degree $U-1$, we write a class by its edge coordinates, and then
compute the norm of each boundary class $\partial_{i*}z_i$.

\begin{proposition}\label{prop:edge-norms}
Let $R$ be $p$-torsionfree, $0\le i<h$, $U\in RO(H)$ be natural,
and $z\in\pi_U^H X_R^{t_i}$. Then
\begin{equation}\label{eq:full-boundary-norm}
N_H^G(\partial_{i*}z)
=\partial_{i*}\bigl(a_H^GN_H^G(z)\bigr).
\end{equation}
The norm is additive on $\pi_{U-1}^H X_R$.
For $U=\rho_H(\mathbf a)$, $\mathbf a\in D_h$, and $z_i\in T_i(U)$,
\[
N_H^G([z_0,\ldots,z_{h-1}])
=\sum_{i=0}^{h-1}
\partial_{i*}\bigl(a_H^GN_H^G(z_i)\bigr).
\]
\end{proposition}

\begin{proof}
The map of nonunital algebras $\partial_i$ in
Proposition~\ref{prop:boundary-indexed-multiplications}
gives \eqref{eq:full-boundary-norm}.

For $[G:H]=p$ and $x,y\in\pi_{U-1}^H X_R$, every term other
than $N_H^G(x)$ and $N_H^G(y)$ in
\cite[equation~(8.34)]{ZhangHZ} contains a product of
$p$ shifted classes, which is zero by \eqref{eq:path-cup}.
Hence
\[
N_H^G(x+y)=N_H^G(x)+N_H^G(y).
\]
Induction preserves shifted degrees, so norm transitivity
gives additivity for all $H<G$.
The formula for $N_H^G([z_0,\ldots,z_{h-1}])$ follows from
\eqref{eq:full-boundary-norm}
and
\[
[z_0,\ldots,z_{h-1}]
=\sum_{i=0}^{h-1}\partial_{i*}z_i.\qedhere
\]
\end{proof}

Let $R$ be $p$-torsionfree and $U=\rho_H(\mathbf a)$, $\mathbf a\in D_h$.
Use Proposition~\ref{lem:period-repair-norm} for $\mathbf a\in D_h^{(0)}$.
For $\mathbf a\in D_h^{(1)}$, the calculation is as follows.
For $U\notin RO_{\mathrm{sel}}(H)$, see Section~\ref{subsubsec:arbitrary-degree-norms}.

At $p=2$, let $h\ge1$ and put $m=2^{n-h}$. We have
\[
\Ind_H^G\sigma_{2^h}
\cong\bigoplus_{a=0}^{m/2-1}
\lambda_{2^n}((2a+1)2^{h-1}).
\]

Let $0\le i<h$ and $z\in\pi_U^H X_R^{t_i}$.
If $\epsilon(\mathbf a)=1$, the cofiber sequence
$
(H/C_{2^{h-1}})_+\longrightarrow S^0
\xrightarrow{a_{\sigma_{2^h}}}S^{\sigma_{2^h}}
$
gives
\[
a_{\sigma_{2^h}}:\pi_{U+\sigma_{2^h}}^H X_R^{t_i}
\twoheadrightarrow\pi_U^H X_R^{t_i}.
\]
Choose $y$ with $a_{\sigma_{2^h}}y=z$. Since
$U+\sigma_{2^h}=\rho_H(a_0+\tfrac12,\ldots,a_{h-1}+\tfrac12,a_h)
\in\rho_H(D_h^{(0)})$,
Proposition~\ref{lem:period-repair-norm} applies to $N_H^G(y)$, and
\[
N_H^G(z)
=N_H^G(a_{\sigma_{2^h}}y)
=a_{\Ind_H^G\sigma_{2^h}}N_H^G(y).
\]

This determines $N_H^G(z)$ for $\mathbf a\in D_h^{(1)}$. For $H=e$ and shifted $V$, $\pi_V^eX_R=0$.

\subsubsection{Norms in arbitrary representation degrees}
\label{subsubsec:arbitrary-degree-norms}
We now extend the norms from $RO_{\mathrm{sel}}(H)$ to $RO(H)$.
Let $H=C_{p^h}<G=C_{p^n}$ and write $V=V_0+D$ as in
Definition~\ref{def:character-translation}.
For $x\in\pi_V^H X_R$, multiplicativity gives, as in the proof of
\cite[Proposition~8.29]{ZhangHZ},
\[
N_H^G(x)=N_H^G(v_D)N_H^G(v_D^{-1}x),
\qquad v_D^{-1}x\in\pi_{V_0}^H X_R.
\]

Complex induction gives
\begin{equation*}
\Ind_H^G\lambda_{p^h}(k)
\cong\bigoplus_{a=0}^{p^{n-h}-1}
\lambda_{p^n}(k+ap^h).
\end{equation*}


For $b(t)=\min(t,p^n-t)$, the standard real weights give
\[
\widetilde v_{\Ind_H^G D}
=\prod_{r,j}\prod_{a=0}^{p^{n-h}-1}
\left(
\frac{\widetilde v_{b(rp^j+ap^h)/p^j,j}}
     {\widetilde v_{b(p^j+ap^h)/p^j,j}}
\right)^{m_{r,j}}.
\]

\begin{proposition}\label{lem:ordered-sign-units}
Let $R$ be $p$-torsionfree. Put
\[
q_l=
\frac{
\deg\bigl(\widetilde v_D^{C_{p^{\min(l,h)}}}\bigr)^{
p^{n-\max(l,h)}}}
{\deg\bigl(\widetilde v_{\Ind_H^G D}^{C_{p^l}}\bigr)}
\qquad(0\le l\le n).
\]
Then
\begin{equation*}
N_H^G(v_D)v_{\Ind_H^G D}^{-1}
=
(\widetilde\theta_{n+1}\varphi)^{-1}
\left(
q_n+\sum_{j=0}^{n-1}
\frac{q_j-q_{j+1}}{p^{n-j}}V^{n-j}(1)
\right).
\end{equation*}
\end{proposition}

\begin{proof}
Let $A(G)$ be the Burnside ring and put
\[
u=N_H^G(\widetilde v_D)\,
\widetilde v_{\Ind_H^G D}^{-1}\in A(G)_{(p)}^\times.
\]
Attention that for $k=rp^j,p^j$, folding the weight $k+ap^h$ requires conjugation
when $a\ge\lceil p^{n-h}/2\rceil$. For $l\le j$ the two signs
cancel. For $l>j$ both fixed spaces vanish. Thus the $C_{p^l}$-mark
of $u$ is $q_l$, and
\[
u=q_n[G/G]+\sum_{j=0}^{n-1}
\frac{q_j-q_{j+1}}{p^{n-j}}[G/C_{p^j}].
\]
Apply the unit map $A(G)_{(p)}\to\pi_0^G X_R$, which sends
\[
[G/C_{p^j}]\longmapsto\tr_{C_{p^j}}^G(1)
=(\widetilde\theta_{n+1}\varphi)^{-1}\bigl(V^{n-j}(1)\bigr).\qedhere
\]
\end{proof}

\subsection{Cyclotomic restriction and norms}\label{subsec:cyclotomic-maps}

We compute $R_{\mathrm{cyc}}$ on the kernel and cokernel in
\eqref{eq:uniform-groups}.
For perfect prisms, the degree-zero norm and cyclotomic restriction
are described in \cite[Section~3.1]{SulymaPrisms}. We start with the following definition.

\begin{definition}\label{def:cyclotomic-coefficients}
Let $R$ be as in Definition~\ref{def:natural-degree}, let $n\ge1$,
and let $U\in RO_{\mathrm{sel}}(C_{p^n})$ be natural. Put
\[
W=\varrho_p^*(U^{C_p}).
\]
For $p$-torsionfree $R$, define
\[
\begin{aligned}
\varphi^{-1}:A_{i+1}(U)\cong\varphi^*A_i(W)
&\longrightarrow A_i(W)
&& (0\le i<n),\\
\varphi^{-1}:T_{i+1}(U)\cong\varphi^*T_i(W)
&\longrightarrow T_i(W)
&& (0\le i<n-1),\\
a\otimes q&\longmapsto\varphi^{-1}(a)q,
\end{aligned}
\]
where the identifications use
\[
\varphi^*(\varphi^{(n-1)*}B)\cong\varphi^{n*}B,
\qquad
\varphi^*I_j\cong I_{j+1}.
\]
\end{definition}

We now describe cyclotomic restriction on vertex and edge coordinates
and compute its iterates on classes of the form $N_H^G(x)$.

\begin{proposition}\label{prop:cyclotomic-operations}
With $R,n,U,W$ as above, cyclotomic restriction is
\[
\begin{aligned}
R_{\mathrm{cyc}}:\ker\Delta_U
&\longrightarrow\ker\Delta_W,\\
(x_0,\ldots,x_n)
&\longmapsto
(\varphi^{-1}x_1,\ldots,\varphi^{-1}x_n),\\[4pt]
R_{\mathrm{cyc}}:\coker\Delta_U
&\longrightarrow\coker\Delta_W,\\
[z_0,\ldots,z_{n-1}]
&\longmapsto
[\varphi^{-1}z_1,\ldots,\varphi^{-1}z_{n-1}],
\end{aligned}
\]
under \eqref{eq:uniform-groups}.
The second map is zero for $n=1$.

For $V=V_0+D$ as in
Definition~\ref{def:character-translation} and
$x_0\in\pi_{V_0}^{C_{p^n}}X_R$,
\[
R_{\mathrm{cyc}}(v_Dx_0)
=
\left(\prod_r r^{m_{r,0}}\right)
v_{\varrho_p^*(D^{C_p})}R_{\mathrm{cyc}}(x_0).
\]
For $p=2$,
\[
R_{\mathrm{cyc}}(a_{\sigma_{2^n}})
=
\begin{cases}
a_{\sigma_{2^{n-1}}},&n\ge2,\\
1,&n=1.
\end{cases}
\]

For $H=C_{p^h}\le C_{p^n}$, $V\in RO(H)$,
$x\in\pi_V^H X_R$, and $0\le a\le n$,
\begin{equation}\label{eq:all-cyclotomic-crossings}
R_{\mathrm{cyc}}^aN_H^{C_{p^n}}(x)
=
\begin{cases}
N_{C_{p^{h-a}}}^{C_{p^{n-a}}}
R_{\mathrm{cyc}}^a(x),&a\le h,\\
N_e^{C_{p^{n-a}}}\psi^{p^{a-h}}
R_{\mathrm{cyc}}^h(x),&a>h.
\end{cases}
\end{equation}
Use the orbit diagonals and $\Phi^{C_{p^a}}\eta$ on the target.
\end{proposition}

\begin{proof}
Put $Y=S^{-U}\wedge X_R$.
Cyclotomic restriction is induced by
\[
Y^{C_{p^n}}
\longrightarrow
(Y^{\Phi_1})^{C_{p^n}}
\simeq (\Phi^{C_p}Y)^{C_{p^n}/C_p}
\xrightarrow[\mathrm{cyclotomic}]{\simeq}
(S^{-W}\wedge X_R)^{C_{p^{n-1}}}.
\]
By \eqref{eq:tate-compression}, $Y\to Y^{\Phi_1}$
projects to the diagram with $Y^{h_0}$ and $Y^{t_0}$ removed.

Assume $R$ is $p$-torsionfree. Since
$W^{C_{p^i}}\cong U^{C_{p^{i+1}}}$,
Proposition~\ref{prop:app-finite-comparison} and
\eqref{eq:app-sign-local-kernels} give
\[
\begin{aligned}
\varphi^{-(i+1)*}A_{i+1}(U)
&\cong
\pi_0\bigl((S^{-W^{C_{p^i}}}\wedge X_R)
^{hC_{p^{n-i-1}}}\bigr)
\cong\varphi^{-i*}A_i(W),
&&0\le i<n,\\
\varphi^{-(i+1)*}T_{i+1}(U)
&\cong
\pi_0\bigl((S^{-W^{C_{p^i}}}\wedge X_R)
^{tC_{p^{n-i-1}}}\bigr)
\cong\varphi^{-i*}T_i(W),
&&0\le i<n-1.
\end{aligned}
\]
The map $Y\to Y^{\Phi_1}$ induces the identity on each middle group.
By Subsection~\ref{subsec:period-line} and
\eqref{eq:app-moving-arrow}--\eqref{eq:app-frob-denominator},
for $0\le i<n-1$,
\[
\begin{aligned}
\ell_i\circ\varphi^{-1}
 &=\varphi^{-1}\circ\ell_{i+1}
 &&:A_{i+1}(U)\longrightarrow T_i(W),\\
\rho_i\circ\varphi^{-1}
 &=\varphi^{-1}\circ\rho_{i+1}
 &&:A_{i+2}(U)\longrightarrow T_i(W).
\end{aligned}
\]
For $\epsilon(\mathbf a)=1$, restrict to the kernels in
\eqref{eq:app-sign-local-kernels} by naturality of the
sign Euler cofiber sequence.
Thus we have a commutative diagram of abelian groups
\[
\begin{tikzcd}[column sep=large,row sep=large]
C_{C_{p^n}}^0(U)
  \arrow[r,"\Delta_U"]
  \arrow[d,"{(x_i)\mapsto(\varphi^{-1}x_{i+1})}"']
& C_{C_{p^n}}^1(U)
  \arrow[d,"{(z_i)\mapsto(\varphi^{-1}z_{i+1})}"] \\
C_{C_{p^{n-1}}}^0(W)
  \arrow[r,"\Delta_W"']
& C_{C_{p^{n-1}}}^1(W).
\end{tikzcd}
\]
Naturality of the homotopy fiber sequence gives the stated
maps on the kernel and cokernel under \eqref{eq:uniform-groups}.

For $0\le a\le h$, \cite[Theorem~2.35]{NormTHH} gives
\[
\Phi^{C_{p^a}}N_{C_{p^h}}^{C_{p^n}}
\simeq
N_{C_{p^{h-a}}}^{C_{p^{n-a}}}\Phi^{C_{p^a}}.
\]
Together with the norm compatibility of the cyclotomic equivalence,
this gives the first case of \eqref{eq:all-cyclotomic-crossings}.
By \cite[Remark~5.2]{SulymaNew} and norm transitivity,
\[
R_{\mathrm{cyc}}N_e^{C_{p^m}}
=
N_e^{C_{p^{m-1}}}R_{\mathrm{cyc}}N_e^{C_p}
=
N_e^{C_{p^{m-1}}}\psi^p
\qquad(m\ge1).
\]
\end{proof}
